\chapter{A \textit{Gekko} szoftver komponens}

Bár a~\ref{sec:ch2-summary} fejezetben azt ígértem, hogy az adatiránynak megfelelő sorrendben mutatom be szoftver
komponenseket, a \textit{Gekko} komponens pedig egy opcionális nulladik lépés, ennek ellenére az utolsó komponens
fejezetben ismertetem.

Ennek legfőbb okai, hogy a \textit{Gekko} komponens használata opcionális, és a~\ref{chp:reporter} fejezetben megismert
szoftverkomponenshez hasonlóan egy külön parancssori futtatható.

A \textit{Gekko} komponens célja, hogy egy megfelelő teszt generátor alkalmazást felhasználva, létrehozni egy olyan
testcsomagot amely kompatibilis a teszt keretrendszerrel. Gyakorlatilag egy \textit{frontendet} biztosít különböző
teszt generátorokhoz, melyek közül tartalmaz egy implementációt a \textit{Csmith}~\cite{csmith} teszt generátorhoz.

Az eredeti ötlet egy egyszerű szkript megvalósítása volt, amely generál valamennyi tesztesetet
\textit{Csmith-el}~\cite{csmith}, majd kimenetüket átkonvertálja a teszt keretrendszer által ismertre. Ennél azonban
generikusabb és robusztusabb megoldást választottam.

A generátor alkalmazás által létrehozott teszteseteket megvizsgálja egy vagy több ismert stabil fordítóprogrammal, hogy
keletkezzen egy viszonyítási alap. Erre azért van szükség mert előfordulhat, hogy egy generált teszteset bizonyos
fordítóprogramon lefordul másikon pedig nem. Az ilyen tesztek nem várt bukásokat eredményezhetnek a tesztelés alatt
álló fordítóprogramot használva. Az ilyen típusú hibákat nehéz megtalálni így jobb elkerülni.

A \textit{Gekko} komponens szoftver architekturális szempontból fel van készítve arra, hogy egyéb teszt generátorokat
is meg tudjon hajtani, a jelen implementáción túl. Ezt többek között a \textit{Python} által nyújtott \textit{Protocol}
osztályokkal valósítottam meg, amik lehetővé teszik a \textit{Struktúrált típusosság} alkalmazását a \textit{Pythonban}
megszokott \textit{Duck típusosság} helyett. Ennek használatával inheritencia nélkül, implicit módon lehet
implementálni interfészeket, a típus korrektség megőrzése mellett.

\newpage

\section{A komponens konfigurációja}

    A \textit{Gekko} szoftver komponens újrahasználja a~\ref{chp:config-manager} fejezetben megismert
    \textit{ConfigManager} komponenst, és definiálja a saját konfigurációs fájl formátumát.

    A konfigurációs fájl az alábbi szekciókat támogatja:

    \begin{itemize}
        \item \textit{DEFAULT}: itt adhatóak meg \textit{Gekko} konfigurációi. A \verb|COMPILE_TIMEOUT| és a
              \verb|RUN_TIMEOUT| amik a verifikáló fordítóprogramok futás idejét, illetve a generált tesztek
              futásidejét korlátozzák. A \verb|PREFIX| kapcsolóval a generált tesztesetek nevének előtagja
              változtatható. A \verb|HOST_COMPILERS| pedig egy lista azokból a fordítóprogramokból amik a viszonyítási
              alapot fogják adni.
        \item \textit{CSMITH}: A jelenleg támogatott generátor kapcsolói. Itt adható meg az \verb|OPTIONS| kapcsolóval
              a program parancssori argumentumai. A \verb|TIMEOUT| kapcsolóval megszorítható a generátor futásideje egy
              időkeretbe.
        \item A további szekciókban a verifikáló fordítóprogramokat lehet felsorolni a \textit{DEFAULT} szekció
              \verb|HOST_COMPILERS| kapcsolója számára. Ezen szekciók kapcsolói, a \verb|COMPILER| ami a program
              elérési útvonala, a \verb|COMPILE_FLAGS| ami a fordítási kapcsolók megadására szolgál, és a
              \verb|LINK_FLAGS| ahol a szerkesztő program számára lehet kapcsolókat átadni.
    \end{itemize}

    A mellékletekben található, \textbf{\ref{fig:gekko-config}. ábrán} látható egy példa a konfigurációs fájlra.

\section{A komponens felépítése}

    A \textit{workers} parancssori argumentummal finomhangolható, az egyidejűleg generált tesztesetek száma. Ez a
    kapcsoló szemantikailag nem összekeverendő, az összesen generálni kívánt tesztesetek számát szabályzó
    \textit{num-tests} kapcsolóval.

    A komponens a teszt keretrendszerhez hasonlóan egy aszinkron esemény hurokban fut. Az esemény hurkon egyszerre
    maximum a \textit{workers} kapcsoló által megadott mennyiségű egyidejű feladat fut. Ezt egy \textit{Szemafor}
    szabályozza amely csak a \textit{workers} parancssori argumentumnak megfelelő feladatot enged egyidejűleg futtatni.
    A \textit{Gekko} komponens viszont egyszerre dupla annyi feladatot tesz fel az esemény hurokra mint amennyi
    egyidejűleg futhat.

    Erre azért van szükség, mert annyi teszteset generálásért felelős feladat jön létre, \\ amennyi a
    \textit{num-tests} értéke. Tehát a feladatok száma lineárisan nő, a generálandó tesztesetek számával. A korlátozás
    azért van, hogy a feladatok számából következő lineárisan növekvő tárigény, felülről korlátos legyen. Továbbá, ha
    gyorsabban regisztrálódnak feladatok az esemény hurkon mint ahogyan elfogynak, akkor a \textit{Szemaforért} is
    zárolási verseny alakulna ki, ezt az esetet is kizárja a felső korlát beiktatása.

    A \textit{Gekko} által futtatott aszinkron feladatok három alkomponensből állnak össze. Ezen komponensek
    \textit{UML Komponens diagramját} ábrázolja a \textbf{\ref{fig:gekko-uml}. ábra}.

    \begin{figure}[ht]
        \centering
        \hrule
        \begin{tikzpicture}
            \begin{umlcomponent}{Gekko}
                \umlbasiccomponent[y=3.3]{Transformer}
                \umlbasiccomponent[x=2]{EvaluationEngine}
                \umlbasiccomponent[x=4,y=3.3]{TestGenerator}

                \umlport{TestGenerator}{0}
                \umlrequiredinterface[interface=Config, with port]{TestGenerator}
                \umlprovidedinterface[interface=SSCTest, with port]{Transformer}

                \umlHVassemblyconnector[interface=TestCase, with port, anchor1=0, anchor2=-60, second arm]{EvaluationEngine}{TestGenerator}
                \umlVHassemblyconnector[interface=GoodTest, with port, anchor1=-120, anchor2=180, first arm]{Transformer}{EvaluationEngine}
            \end{umlcomponent}

            \umlHVHrelation{Gekko-east-port}{TestGenerator-east-interface}
            \umlHVHrelation{Transformer-west-interface}{Gekko-west-port}

            \umlport{Gekko}{0}
            \umlport{Gekko}{180}
        \end{tikzpicture}
        \hrule
        \caption{\label{fig:gekko-uml} \textit{Gekko UML Komponens} diagram}
    \end{figure}

    \newpage

    Az architektúra gyakorlatilag egy csővezeték felépítést követ. Fontos kiemelni, hogy ha a csővezetékben bárhol hiba
    keletkezik akkor a rendszer megpróbálja a folyamatot a csővezeték elejéről. Ezt maximum hatszor teszi meg, ha ennyi
    alkalommal nem sikerül tesztet generálni, akkor feltételezhetjük hogy egyéb, az alkalmazás működésén kívül eső
    probléma áll fent. Az újrapróbálkozás felső korlátja azt is jelenti, hogy a kért teszteseteknél előfordulhat, hogy
    kevesebbet generál az alkalmazás. A bementi konfiguráció alapján a \textit{Gekko} a következő lépés sorozatot
    hajtja végre, alkomponensek szerint:

    \begin{enumerate}
        \item Teszteset generálás.
        \item Teszteset verifikáció, a konfigurációban definiált fordítóprogramokkal.
        \item Verifikált teszteset transzformáció, a \textit{SuperSet} teszt keretrendszer formátumára.
    \end{enumerate}

    A komponensek számára a konfigurációt elérhetővé kell tenni egy globális helyen, ezt a feladatot
    \textbf{\ref{fig:gekko-uml}. ábrán} nem szereplő, \textit{Collector} komponens látja el.

    A \textit{Collector} komponens az alábbi fontos konfigurációkat teszi elérhetővé a konfigurációs fájlból
    beolvasottakon túl:

    \begin{itemize}
        \item A teszteset generáló objektumot, amelyen keresztül egy közös interfészen keresztül hívhatóak különböző
              teszteset generátorok.
        \item A verifikáló fordítóprogram objektumokat.
    \end{itemize}

    A komponens megvizsgálja a verifikáló fordítóprogramokat pár egyszerű invokációval, így megbizonyosodik róla hogy a
    parancssori argumentum interfészük a \textit{Gekko} által ismert. A komponens ezen felül több a futás során
    szükséges elérési útvonal létrehozásáról is gondoskodik, például a kimeneti mappáéról, és a \textit{log} fájlok
    mappáéról.

\section{A \textit{Gekko} alkomponensei}

    Ahogyan azt a \textbf{\ref{fig:gekko-uml}. ábrán} láthattuk az első komponens a \textit{TestGenerator}.

    A \textit{TestGenerator} alkomponens az egyetlen publikus tagfüggvényén, a \textit{generate} függvényen keresztül
    generál egy tesztesetet. Ezt a konfigurációban megadott, teszteset generátor alkalmazás alprocesszusban történő
    meghívásával teszi.

    A második komponens az \textit{EvaluationEngine}, vagyis a kiértékelő komponens. Ez a komponens végzi a
    \textit{TestGenerator} által generált tesztesetek verifikációját. A konfigurációban megadott fordítóprogramokkal
    lefordítja a tesztesetet.

    Azonban itt nem fejezi be a munkát. A tény, hogy minden alkalmazott fordítóprogram tudott futtatható állományt
    generálni, még nem elég ahhoz, hogy megbizonyosodjunk a teszteset helyességéről.

    A \textit{Csmith}~\cite{csmith} teszteset generáló program kimeneti forrásfájljai futásidejű tesztesetek.
    Kimenetükben egy ellenőrzőösszeget adnak vissza. Ennek az összegnek a helyességét kell verifikálni. Ennek
    tudatában, az \textit{EvaluationEngine} implementációja az összes tesztet le is futtatja, és az eredményüket
    ellenőrzi, hogy az elvárt ellenőrzőösszeget produkálják.

    Eltérés esetén a teszteset nem kerül felhasználásra és a csővezeték hibával megáll. Itt a fentebb említett módon
    történhet új tesztesettel újra próbálkozás.

    A harmadik alkomponens a \textit{Transformer}. Ahogyan azt az \textit{EvaluationEngine} kapcsán említettem a
    \textit{Csmith}~\cite{csmith} kimenete egy specifikus formátumban van.

    A formátum tartalmaz metainformációkat a generált tesztesetről, és a generáló programról. Ezek nagyon hasznos
    információk, mert nem kell a tesztesetet lefordított futtatható állományát elmenteni, ha reprodukálni szeretnénk a
    bukást. Egyszerűen a metainformációból újra előállítható a teszteset. A \textit{Transformer} a következő
    metainformációkat tárolja el:

    \begin{itemize}
        \item A generáló alkalmazást, verziójával együtt.
        \item A generáló alkalmazás parancssori argumentumait.
        \item A \textit{Gekko} saját egyedi azonosítóját \textit{(hash)}.
        \item A tesztgeneráló által alkalmazott \textit{seed} érték, amely a véletlen generálást reprodukálhatóvá
              teszi.
    \end{itemize}

    A metainformációkat egy \textit{JSON}~\cite{json} fájlban tárolja a komponens. Az alábbi
    \textbf{\ref{fig:metadata-json}. ábrán} látható módon.

    \newpage

    \begin{figure}[ht]
        \centering
        \begin{minted}{json}
{
    "test_1.c": {
        "generator": "csmith 2.4.0",
        "options": "--no-argc --max-array-dim 5 --no-longlong --no-packed-struct -o test_1.c",
        "revision": "0ec6f1b",
        "seed": "15774984650604141557"
    }
}
        \end{minted}
        \caption{\label{fig:metadata-json} Egy teszteset metainformációja \textit{JSON} formátumban}
    \end{figure}

    \subsection{Az alkomponens interfészek}

        A különböző alkomponensek egy-egy interfész rétegen keresztül kommunikálnak egymással. Ez az interfész réteg a
        \textit{Python} protokollokra építkezik.

        Ahogyan azt a fejezet elején is említettem a protokollokkal elérhető a \textit{struktúrált típusosság}
        öröklődés beveztése nélkül. Ez nagyban leegyszerűsíti a kódot, ugyanis elhagyhatók a \textit{Pythonban}
        megszokott \verb|@abstractmethod| függvény annotációk.

        Ezek elkerülése nem csak a kódolvashatóságon javít, hanem a futás időn is. Ez a javaulás elenyésző ugyan de
        attól még említésre méltó. A javulás annak köszönhető, hogy az \textit{abstract} annotációkat a
        \textit{Python}, mint ahogyan minden más konstrukciót is, futás időben értékeli ki.

        Ezzel szemben, a protokollokat a \textit{Python interpreter} futásidőben nem ellenőrzi. Az ellenőrzésükhöz
        szükség van egy harmadik féltől származó típus ellenőrző alkalmazásra. Például, \textit{mypy}~\cite{mypy},
        \textit{pyright}~\cite{pyright}, \textit{ty}~\cite{ty}.

        A protokollok egy szoftverarchitekturális szempontból, is előnyesek, mert egy tisztább megoldást kínálnak mint
        az egyébként \textit{Pythonban} elérhető \textit{abstract} osztályok. Mivel nem osztályok és nincs szükség
        öröklődésre, így adattagokat sosem kapnak az implementációk a protokolloktól.

        A protokollok hátránya egyedül a lefelé kasztolásnál jön elő, mert csak a \\ \verb|@runtime_checkable|
        annotációval ellátott protokollok tesztelhetők futásidőben. Erre a felhasználásra ideálisabb \textit{abstract}
        osztályokat alkalmazni.
